\chapter{Diffusion models and image generation}
\label{ch:diffusion}

\begin{quote}
\emph{``Diffusion models generate images by learning to reverse a noise process. Start with pure noise, denoise step by step, and a coherent image emerges.''}
\end{quote}

% ============================================================
\section{The diffusion process}

\subsection{Forward process (adding noise)}

Given an image $\mathbf{x}_0$, the forward process adds Gaussian noise over $T$ timesteps:
\[
q(\mathbf{x}_t | \mathbf{x}_{t-1}) = \mathcal{N}(\mathbf{x}_t; \sqrt{1-\beta_t}\,\mathbf{x}_{t-1},\; \beta_t \mathbf{I})
\]

where $\beta_t$ is the noise schedule. After enough steps, $\mathbf{x}_T \approx \mathcal{N}(\mathbf{0}, \mathbf{I})$.

\subsection{Reverse process (denoising)}

A neural network learns to reverse the noise:
\[
p_\theta(\mathbf{x}_{t-1} | \mathbf{x}_t) = \mathcal{N}(\mathbf{x}_{t-1}; \boldsymbol{\mu}_\theta(\mathbf{x}_t, t),\; \sigma_t^2 \mathbf{I})
\]

The model predicts the noise $\boldsymbol{\epsilon}_\theta(\mathbf{x}_t, t)$ and uses it to compute $\boldsymbol{\mu}_\theta$.

\subsection{Training objective}

The simplified DDPM loss is:
\[
L = \mathbb{E}_{t, \mathbf{x}_0, \boldsymbol{\epsilon}} \left[\| \boldsymbol{\epsilon} - \boldsymbol{\epsilon}_\theta(\mathbf{x}_t, t) \|^2 \right]
\]

\begin{minted}{python}
import torch
import torch.nn.functional as F

def diffusion_loss(model, x_0, noise_scheduler):
    """Simplified DDPM training step."""
    batch_size = x_0.shape[0]
    # Sample random timesteps
    t = torch.randint(0, noise_scheduler.num_train_timesteps,
                      (batch_size,), device=x_0.device)
    # Sample noise
    noise = torch.randn_like(x_0)
    # Add noise to images
    x_t = noise_scheduler.add_noise(x_0, noise, t)
    # Predict noise
    noise_pred = model(x_t, t).sample
    # MSE loss
    loss = F.mse_loss(noise_pred, noise)
    return loss
\end{minted}

% ============================================================
\section{Noise schedules}

The noise schedule $\{\beta_t\}_{t=1}^T$ controls how fast noise is added:

\begin{itemize}
  \item \textbf{Linear:} $\beta_t$ increases linearly from $\beta_1 = 10^{-4}$ to $\beta_T = 0.02$.
  \item \textbf{Cosine:} smoother noise addition, better for small images.
  \item \textbf{Scaled linear:} used by Stable Diffusion, tuned for latent space.
\end{itemize}

\begin{minted}{python}
import numpy as np
import matplotlib.pyplot as plt

T = 1000
# Linear schedule
beta_linear = np.linspace(1e-4, 0.02, T)
alpha_linear = np.cumprod(1 - beta_linear)

# Cosine schedule
s = 0.008
steps = np.arange(T + 1) / T
f = np.cos((steps + s) / (1 + s) * np.pi / 2) ** 2
alpha_cosine = f[1:] / f[0]

plt.figure(figsize=(8, 4))
plt.plot(alpha_linear, label="Linear")
plt.plot(alpha_cosine, label="Cosine")
plt.xlabel("Timestep")
plt.ylabel("Cumulative alpha (signal remaining)")
plt.title("Noise schedules")
plt.legend()
plt.tight_layout()
plt.savefig("noise_schedules.png", dpi=150)
plt.show()
\end{minted}

% ============================================================
\section{The U-Net architecture}

Diffusion models typically use a \textbf{U-Net}: an encoder-decoder with skip connections. The U-Net takes a noisy image and the timestep, and predicts the noise.

Key components:
\begin{itemize}
  \item \textbf{Downsampling blocks}: convolutions that reduce spatial resolution.
  \item \textbf{Upsampling blocks}: transposed convolutions that increase resolution.
  \item \textbf{Skip connections}: concatenate encoder features to decoder features.
  \item \textbf{Timestep embedding}: sinusoidal embedding of $t$, injected into each block.
  \item \textbf{Cross-attention}: for text-conditioned generation, the U-Net attends to text embeddings.
\end{itemize}

% ============================================================
\section{Stable Diffusion}

Stable Diffusion operates in \emph{latent space}: images are compressed by a VAE encoder, diffusion happens in the latent space, and the VAE decoder reconstructs the image.

\begin{minted}{python}
from diffusers import StableDiffusionPipeline
import torch

pipe = StableDiffusionPipeline.from_pretrained(
    "stabilityai/stable-diffusion-2-1-base",
    torch_dtype=torch.float16,
)
pipe = pipe.to("cuda")

prompt = "A serene African village at sunset, digital art, highly detailed"
image = pipe(
    prompt,
    num_inference_steps=30,
    guidance_scale=7.5,
).images[0]

image.save("african_village.png")
image
\end{minted}

\begin{aitip}
\texttt{guidance\_scale} controls how closely the image follows the prompt. Values of 7--9 work well for most prompts. Higher values produce more literal interpretations but can reduce image quality.
\end{aitip}

% ============================================================
\section{Negative prompts and parameters}

\begin{minted}{python}
image = pipe(
    prompt="Portrait of a scientist in a laboratory, photorealistic",
    negative_prompt="blurry, low quality, deformed, cartoon",
    num_inference_steps=50,
    guidance_scale=8.0,
    height=512,
    width=512,
).images[0]
image.save("scientist.png")
\end{minted}

% ============================================================
\section{Image-to-image generation}

Start from an existing image and modify it:

\begin{minted}{python}
from diffusers import StableDiffusionImg2ImgPipeline
from PIL import Image

pipe_img2img = StableDiffusionImg2ImgPipeline.from_pretrained(
    "stabilityai/stable-diffusion-2-1-base",
    torch_dtype=torch.float16,
).to("cuda")

init_image = Image.open("sketch.png").resize((512, 512))

result = pipe_img2img(
    prompt="A beautiful watercolor painting of a landscape",
    image=init_image,
    strength=0.75,  # 0=no change, 1=complete regeneration
    guidance_scale=7.5,
).images[0]
result.save("watercolor.png")
\end{minted}

% ============================================================
\section{ControlNet: guided generation}

ControlNet adds spatial conditioning (edges, depth, pose) to Stable Diffusion:

\begin{minted}{python}
from diffusers import StableDiffusionControlNetPipeline, ControlNetModel
from controlnet_aux import CannyDetector
from PIL import Image
import torch

# Load ControlNet for Canny edges
controlnet = ControlNetModel.from_pretrained(
    "lllyasviel/sd-controlnet-canny", torch_dtype=torch.float16
)
pipe_cn = StableDiffusionControlNetPipeline.from_pretrained(
    "stabilityai/stable-diffusion-2-1-base",
    controlnet=controlnet,
    torch_dtype=torch.float16,
).to("cuda")

# Extract edges from reference image
canny = CannyDetector()
reference = Image.open("building_photo.png")
edges = canny(reference, low_threshold=100, high_threshold=200)

# Generate with edge guidance
result = pipe_cn(
    prompt="A futuristic building, cyberpunk style, neon lights",
    image=edges,
    num_inference_steps=30,
).images[0]
result.save("cyberpunk_building.png")
\end{minted}

% ============================================================
\section{Batch generation and seed control}

\begin{minted}{python}
import torch

# Reproducible generation with seeds
generator = torch.Generator("cuda").manual_seed(42)

images = pipe(
    prompt=["A cat astronaut", "A dog scientist", "A bird musician"],
    num_inference_steps=30,
    guidance_scale=7.5,
    generator=generator,
).images

for i, img in enumerate(images):
    img.save(f"generated_{i}.png")
\end{minted}

\begin{exercise}
\begin{enumerate}
  \item Generate 5 images with the same prompt but different seeds. How much variation do you observe?
  \item Experiment with \texttt{guidance\_scale} values: 1, 5, 7.5, 12, 20. Document the effect on image quality and prompt adherence.
  \item Use image-to-image to transform a simple sketch into a detailed illustration. Try different \texttt{strength} values.
  \item Apply ControlNet with Canny edge detection to generate architectural variations of a building photo.
  \item Generate an image, then use img2img to create 3 variations with increasing \texttt{strength} (0.3, 0.5, 0.8).
\end{enumerate}
\end{exercise}

\begin{ethicsbox}
Image generation raises serious ethical concerns:
\begin{itemize}
  \item \textbf{Deepfakes:} generating realistic images of real people without consent.
  \item \textbf{Copyright:} models trained on copyrighted art without artist permission.
  \item \textbf{Bias:} models can perpetuate stereotypes (e.g., generating mostly light-skinned faces for ``professional'').
  \item \textbf{Harmful content:} models can generate violent or explicit images.
\end{itemize}
Always use safety checkers (enabled by default in diffusers), watermark AI-generated images, and never generate non-consensual images of real people.
\end{ethicsbox}

% ============================================================
\section{Chapter summary}

\begin{itemize}
  \item Diffusion models learn to reverse a noise process: start from noise, denoise step by step.
  \item The training objective is simple: predict the noise that was added.
  \item Noise schedules (linear, cosine) control the diffusion speed.
  \item U-Net with cross-attention is the standard architecture for text-conditioned generation.
  \item Stable Diffusion works in latent space for efficiency.
  \item ControlNet adds spatial guidance (edges, depth, pose) to the generation process.
  \item Guidance scale, negative prompts, and number of steps are the key generation parameters.
  \item Ethical deployment requires safety filters, watermarking, and consent.
\end{itemize}
